\chapter*{知识点首次出现索引}
\label{hw:knowledge-index}
\pdfbookmark[chapter]{知识点首次出现索引}{hw-knowledge-index}
\noindent 下表指向每题后的完整知识讲解，不仅指向计算答案。知识在首次出现处展开；后续使用时结合题内“已学前置”回查。题号采用“章号—原题号”，页码可点击跳转。

\begingroup
\setstretch{1.0}
\small
\setlength{\tabcolsep}{5pt}
\begin{longtable}{@{}p{1.05cm}p{12.25cm}r@{}}
\toprule
题号 & 知识点与核心公式范围 & 页码\\
\midrule
\endfirsthead
\toprule
题号 & 知识点与核心公式范围（续） & 页码\\
\midrule
\endhead
\bottomrule
\endfoot
2—1 & 周期、$f$与$\omega$；均值、绝对均值、均方/有效值、方差；峰值与峰峰值、峰值/波形因数；直流和交流均方关系；能量与功率分类。 & \pageref{kn:2-1}\\
2—2 & 复数模、相位、共轭及欧拉公式；傅里叶正反变换、存在条件和谱密度单位；Hz/角频率约定；线性、时移、时标伸缩、共轭对称；能量Parseval关系与指数时间尺度。 & \pageref{kn:2-2}\\
2—3 & 冲激筛选、权重和缩放；主值分布；微分及积分的变换；零频项的确定；常数、复指数、正余弦、阶跃、符号函数的广义变换对。 & \pageref{kn:2-3}\\
2—4 & 频移与时移的区别；卷积定义、冲激卷积及乘积定理；矩形窗、Sa/sinc；主瓣、旁瓣、谱泄漏及观察时长；有符号谱与幅谱。 & \pageref{kn:2-4}\\
2—5 & 单边阻尼振荡的衰减率、包络和振荡周期；复频谱的相位象限；幅谱峰位与名义频率；能量积分、高频尾部及边界光滑性。 & \pageref{kn:2-5}\\
2—6 & 基带/载波/边带/带宽、频带重叠；抑制载波与包络符号；同步解调；傅里叶级数两种形式、正交性、单/双边幅值及周期功率Parseval关系。 & \pageref{kn:2-6}\\
3—1 & LTI与零状态；微分方程、传递函数、冲激响应、卷积和频率响应；相量稳态法；一阶模型、静态增益与时间常数；稳态/暂态及不失真条件。 & \pageref{kn:3-1}\\
3—2 & 幅值误差基准、时间常数不等式；阶跃63.2\%、上升/调节时间；截止频率与精度带宽；Bode幅相渐近；幅值与相位联合设计。 & \pageref{kn:3-2}\\
3—3 & 质量—阻尼—刚度模型；$\omega_n$、$\zeta$与极点；阻尼分类、欠阻尼阶跃/超调/调节时间；二阶幅相与共振峰推导、幅值平坦与低频相位。 & \pageref{kn:3-3}\\
3—4 & 动态灵敏度、静态标定与单频动态校正；按工作频带选择固有频率；相位延迟和群延迟的定义、单位、区别与适用范围。 & \pageref{kn:3-4}\\
4—1 & 应变及电阻定律；应变系数、泊松效应与电阻率变化；方向、温度交叉影响和自热；电流读出、分辨力及电桥读出动机。 & \pageref{kn:4-1}\\
4—2 & 电容三类变参、差动读出；磁阻/自感/互感与差动变压器；复阻抗及交流桥平衡；参数到电压/频率的转换、相敏读出。 & \pageref{kn:4-2}\\
4—3 & 外光电、光导与光伏效应；光谱阈值、响应度和量子效率；光路分类、光电编码与计数；暗信号、饱和和噪声等读出限制。 & \pageref{kn:4-3}\\
4—4 & 洛伦兹力平衡到霍尔电压；霍尔系数、方向与偏置；测磁场/电流/转速的转换链；霍尔式与磁电感应式的静动态区别。 & \pageref{kn:4-4}\\
4—5 & 涡流与趋肤深度、导电性和缺陷取向；压电超声发射/接收、声程、声阻抗与回波；渗透、磁性检测的适用范围和材料限制。 & \pageref{kn:4-5}\\
4—6 & 光纤导光、强度/相位/波长读出；反射式位移传感；膜片压力到光学量的转换；FBG与温度交叉、机械与光电带宽。 & \pageref{kn:4-6}\\
4—7 & 量程与灵敏度；分辨力、准确度、非线性、迟滞、重复性、漂移和交叉敏感；误差基准；加载、后级匹配与完整测量链选型。 & \pageref{kn:4-7}\\
5—1 & 四臂桥路平衡、精确分压与线性化；单/双/全桥的接法、灵敏度和非线性；温度补偿；戴维南等效与负载效应。 & \pageref{kn:5-1}\\
5—2 & 交流应变桥的抑制载波调幅；三种线谱约定的实际换算；同步/相敏解调、参考相位及恢复符号的条件。 & \pageref{kn:5-2}\\
5—3 & 常规AM、调制度及多音包络；边带幅值/权重；同阻抗下的载波与边带功率、效率；调幅带宽与包络检波条件。 & \pageref{kn:5-3}\\
5—4 & 激励载频与被测量工作频带；搬移—放大—解调—低通链；直流基带与极性恢复；各环节通带及整体幅相误差。 & \pageref{kn:5-4}\\
5—5 & 滤波类型、通/阻带、截止点、阶次与滚降；幅值/功率dB；几何中心、分数倍频程、相对带宽和品质因子及约定边界。 & \pageref{kn:5-5}\\
5—6 & 一阶RC实现、负载和缓冲、级联响应；多频幅相比改变；RC相位/群延迟的应用；滤除高频与保持有效信号波形的取舍。 & \pageref{kn:5-6}\\
\end{longtable}
\endgroup
